\label{ch:related-work}

\section{MoE communication and load balancing}

Most MoE literature treats balance at training time, routing time, or expert
placement time. Load-balancing losses and bias updates aim to prevent router
collapse \cite{fedus2021switch,wang2024lossfree}. Distributed inference adds a
systems bottleneck: expert parallelism interleaves computation with all-to-all
communication. Lina shows that skewed expert popularity creates imbalanced
transfer sizes and bandwidth demand, and reports substantial P95 inference
latency improvements from communication-aware scheduling
\cite{li2023lina}. COMET independently reports that inter-device
communication can account for up to 47\% of a MoE layer's execution time in
the evaluated models and frameworks \cite{zhang2025comet}. These measurements
motivate treating network interference as a first-class serving concern. They
do not, however, imply a universal fixed link-utilization ratio.

This work begins after EPLB has already chosen a valid new placement. It does
not change router probabilities, model outputs, or the placement objective.
Instead, it schedules the physical copying needed to realize that placement.
A placement algorithm decides \emph{what} state should move; the migration
scheduler decides \emph{when} independent transfers should share network
resources. The two decisions are complementary.

\section{Collective-flow scheduling}

A set of expert migrations resembles a coflow: several point-to-point flows
share a common completion or interference objective. Varys models coflow
completion through the most heavily loaded ingress or egress endpoint and the
remaining measured bandwidth at that endpoint \cite{chowdhury2014varys}. This
bottleneck view is closely related to the maximum normalized link-load term in
the ILP developed in \cref{ch:formulation}. The present model is intentionally
simpler: it uses fixed routes and a normalized foreground-traffic parameter,
and it optimizes a serving-interference surrogate rather than coflow
completion time.

\section{Graph scheduling and expert migration}

A migration layer can be represented by a graph whose vertices are ranks and
whose directed edges are remote flows. If each rank may participate in at most
one flow in a batch, edges sharing an endpoint cannot receive the same batch
identifier. This is an edge-coloring constraint \cite{west2001graph}. Transfer
direction and expert identity remain necessary for execution, but direction
does not change whether two flows contend at an endpoint. Each generated batch
is a matching in this rank graph.

Exact edge coloring alone is not the complete systems objective: topology,
transfer sizes, replica choices, and foreground traffic also matter. The
deployed first-fit rule is therefore a conservative, topology-independent
heuristic rather than a claim of general optimality. It is inspired by the
migration grouping used in OmniInfer's placement manager
\cite{omniinferplacement}. The contribution of this thesis is the explicit
interference formulation and the integration of this rule into vLLM's
asynchronous EPLB path, including deterministic flow construction, NIXL
lifecycle handling, cycle-wide schedule precomputation, correctness tests,
and end-to-end serving evaluation.
